\section{Later work}
\label{sec:related}

This is a brief orientation, and it relies on the papers' own accounts of
their results.

\paragraph{Before and alongside.}
Ridge~\cite{ridge10} gave a rely-guarantee proof system for x86-TSO,
mechanized in HOL4 and applied to Simpson's four-slot algorithm. Cohen and
Schirmer~\cite{cohen10} and Owens~\cite{owens10} took the other route:
conditions on a program, an ownership discipline and triangular-race
freedom respectively, under which its TSO executions are sequentially
consistent, so that buffers need not be reasoned about at all. The
dissertation, and the proposal with Berdine that preceded
it~\cite{wehrman11}, sat between these: buffers are visible in assertions,
and separation is meant to keep the reasoning local.

\paragraph{iCAP-TSO.}
The closest successor is the logic of Sieczkowski, Svendsen, Birkedal and
Pichon-Pharabod~\cite{sieczkowski15}, which cites the dissertation. It
extends a higher-order concurrent separation logic with store buffers and
with connectives for them, including one stating that an assertion holds
until a buffered update is flushed and another afterwards. That is a
temporal connective in the dissertation's sense, attached to a single
update. Their central idea is a \emph{fiction of sequential consistency}:
a library that synchronizes enough can be given a specification from which
its clients reason in ordinary separation logic. They observe, fairly, that
the dissertation's logic cannot exploit synchronization to simplify client
proofs in that way. They also verify a spin lock whose release is a plain
buffered store, which is exactly the unlocked access to shared state that
Section~\ref{sec:open} leaves open here. Jacobs~\cite{jacobs14} proposed
``TSO spaces'' for separation logic with a related aim.

\paragraph{Temporal reasoning elsewhere.}
Bornat, Alglave and Parkinson~\cite{bornat15} built a rely-guarantee logic
for weak memory around temporal modalities, ``since'' and ``so far'' among
them. The idea that ordering between writes belongs in the assertion
language, which the dissertation expressed as a separating conjunction, has
thus reappeared in other forms.

\paragraph{Other memory models.}
Most later work targets the C11 model and weaker ones rather than TSO:
relaxed separation logic~\cite{vafeiadis13} and GPS~\cite{turon14}, and
their reconstruction inside Iris~\cite{kaiser17}, which made mechanized
soundness proofs the norm for such logics. For x86 itself,
Pierogi~\cite{bila22} is an Owicki--Gries logic over a view-based semantics
that also covers persistent memory.

\paragraph{Where this report stands.}
Against that background, what the repaired logic offers is narrow and, as
far as I know, still distinctive: sequential and spatial separating
conjunctions over the whole history of a thread's writes, each with a frame
rule, and a soundness proof checked end to end against an operational TSO
machine. Section~\ref{sec:trilemma}'s conflict between temporal framing,
units and ownership transfer does not depend on TSO and may be of use to
anyone designing a logic with a sequential conjunction. What it lacks,
unlocked access and a way for libraries to hide the memory model from
their clients, is what the later logics were built to provide.
